\paragraph{The selected tangent degrees.}
All degrees in this paragraph are ordinary Arakelov degrees divided by
the degree of the number field. Let $\bar V$ have rank $n$ and degree
$d$, with pairwise orthogonal, generically complementary integral
sublattices $W_\tau,W'_\tau$. Let $V_\tau\subset W_\tau$ have the same
generic fibre and the inherited metric. Put
\[
 X=\sum_\tau\widehat{\deg}_{\rm n}W_\tau,\quad
 Y=\sum_\tau\widehat{\deg}_{\rm n}W'_\tau,\quad
 T=\sum_\tau\widehat{\deg}_{\rm n}V_\tau,\quad
 R=\sum_\tau\operatorname{rk}V_\tau .
\]
The two integral indices give
\begin{equation}\label{geo-selected-index-identity}
 \delta=d-X-Y\ge0,\qquad \nu=X-T\ge0,\qquad
 b=X-Y,\qquad d=2T-b+\delta+2\nu .
\end{equation}
There is no archimedean correction, since all summands are orthogonal
and retain the given metrics.

\begin{proposition}\label{geo-selected-tangent-sum}
If $C_s\ge0$ and
\[
 T\le RC_s\log(2+n+(-d)_+),
\]
then
\[
 T\le RC_s\log\bigl(2+n+(b-\delta-2\nu)_+\bigr).
\]
If also $\Lambda\ge1$, $C_f\ge0$, $n\le G\le e^\Lambda$ and
$b\le C_fG\Lambda$, then
\begin{equation}\label{geo-selected-log-budget}
 T\le C_A R\Lambda,\qquad C_A=C_s\{1+\log(3+C_f)\}.
\end{equation}
\end{proposition}
\begin{proof}
For $T\le0$ the first conclusion is immediate. For $T>0$,
\eqref{geo-selected-index-identity} gives
$(-d)_+=(b-\delta-2\nu-2T)_+\le(b-\delta-2\nu)_+$.
For the second conclusion, bound the logarithm's argument by
$e^\Lambda(3+C_f\Lambda)$ and use concavity:
\[
 \log(3+C_f\Lambda)
 \le\log(3+C_f)+\frac{C_f}{3+C_f}(\Lambda-1)
 \le\Lambda\log(3+C_f).
\]
The proof includes $R=0$ and $C_s=0$.
\end{proof}

Only the signed total $b$ occurs; both integral defects are favourable.
If the selected eigenspace has dimension $m$, an adapted family has
$r$ toric vectors and $m-r$ abelian vectors, and there are $f$ conjugate
pairs, then the exact selected rank is $R=f(m-r)$. Thus a comparison
\[
 \Phi\le fmX_0+T-fr\log\operatorname{rd}(E)
                         +\frac{fr}{2}\log(4\pi)
\]
together with $X_0\le C_X\Lambda$, $\log\operatorname{rd}(E)\ge0$ and
$\log(4\pi)\le\Lambda$ gives
\[
 \Phi\le\{C_X+\max(C_A,\tfrac12)\}fm\Lambda .
\]
In the purely toric case the selected abelian sum is zero.

\paragraph{A short integral jet argument.}
Write $\operatorname{av}_K$ for the average over all complex embeddings
of $K$, counting conjugate embeddings separately. The next lemma
isolates the elementary determinant argument.

\begin{lemma}\label{geo-leading-jet-determinant}
Let $J_0,\ldots,J_k$ be independent $K$-linear forms on $K^r$ and
$W_i=\bigcap_{j<i}\ker J_j$. Suppose
$J_i(W_i\cap\mathcal O_K^r)\subset\mathcal O_K$ and, for every embedding,
\[
 \|J_i|_{W_i\otimes_\sigma\mathbb C}\|_{\rm op}\le A_{i,\sigma},
 \qquad A_{i,\sigma}>0 .
\]
Then $0\le\sum_{i=0}^k\operatorname{av}_K\log A_{i,\sigma}$.
\end{lemma}
\begin{proof}
At each finite prime, choose an integral basis adapted to the saturated
kernel flag. The row matrix is triangular in its first $k+1$ columns,
zero in the remaining columns, and has integral diagonal entries.
Every coordinate of $J_0\wedge\cdots\wedge J_k$ is therefore integral.
For a nonzero such coordinate $c$, row Gram--Schmidt gives
\[
 |\sigma(c)|\le\|J_0\wedge\cdots\wedge J_k\|_\sigma
       =\prod_{i=0}^k\|J_i|_{W_i\otimes_\sigma\mathbb C}\|_{\rm op}
       \le\prod_{i=0}^kA_{i,\sigma}.
\]
Now use $|N_{K/\mathbb Q}(c)|\ge1$.
\end{proof}

The required integrality has a direct geometric interpretation.
Let $\mathcal G/\mathcal O_K$ be a group scheme, let $\mathcal L$
be a line bundle rigidified at its identity, and let $x$ be an integral
invariant tangent vector. For integral sections of $\mathcal L^{\otimes n}$,
take successive directional derivatives in a generic local frame
with that prescribed value at the identity. At a finite prime choose
an integral frame with the same value. Derivatives of its regular
representing functions are integral. When the earlier derivatives
vanish, the Leibniz rule shows that the leading derivative is unchanged
by the transition function, whose value is one. It is consequently
integral at every finite prime. These are undivided derivatives:
no assertion about division by $i!$ is needed. The same argument
identifies the leading derivative in a holomorphic frame with the
same rigidification.

\begin{proposition}\label{geo-theta-jet-bound}
Suppose the generic fibre is an abelian variety $A/K$ of dimension
$g\ge1$, $L$ is ample, and $x\ne0$ is an integral tangent vector.
Let $n=3k$, $k\ge1$. Suppose $r$ integral sections of $L^{\otimes n}$
span its generic section space. Write $\ell_\sigma=\|x\|_\sigma>0$.
Assume their functions in the rigidified theta frame satisfy
\[
 |f_c(zx)|\le\sqrt r\,M_\sigma\|c\|_2
                    \exp(n\ell_\sigma^2|z|^2/2),
 \qquad M_\sigma>0,\quad c\in\mathbb C^r .
\]
For $U=\operatorname{av}_K\log M_\sigma$ one has
\begin{equation}\label{geo-exact-jet-cost}
 -\operatorname{av}_K\log\ell_\sigma
       \le\frac12\log(enk)+\frac{\log r+2U}{k}.
\end{equation}
\end{proposition}
\begin{proof}
The third power of an ample line bundle on an abelian variety is
very ample \cite[Theorem~2.2.3]{genestier-ngo-shimura}.
Choose $u_0,u_1\in H^0(A,L^{\otimes3})$ with $u_0(0)\ne0$,
$u_1(0)=0$, and $D_x(u_1/u_0)(0)\ne0$. Then
$u_0^{k-i}u_1^i$ has exactly order $i$ along $\exp(zx)$,
for every $0\le i\le k$. The supplied spanning sections therefore
have independent jet rows through order $k$. Only this generic
rank argument uses $u_0,u_1$; their integrality is unnecessary.

For $1\le i\le k$, apply Cauchy's formula at radius
$\sqrt{i/n}/\ell_\sigma$ and use $i!\le i^i$. The operator bound
on the $i$-th leading jet is
\[
 A_{i,\sigma}=\sqrt r\,M_\sigma
                    (\sqrt{enk}\,\ell_\sigma)^i .
\]
The same formula holds for $i=0$. Leading-jet integrality and
Lemma~\ref{geo-leading-jet-determinant} give
\[
 0\le(k+1)(\tfrac12\log r+U)
       +\frac{k(k+1)}2
          \{\tfrac12\log(enk)+\operatorname{av}_K\log\ell_\sigma\}.
\]
Division proves \eqref{geo-exact-jet-cost}.
\end{proof}

The norm convention is essential: for the prescribed Hodge metric,
$\ell_\sigma^2=\pi H_\sigma(x,x)$ and the theta weight is
$\exp(-\pi nH_\sigma(w,w)/2)$. These give exactly the Gaussian
bound above. Also, $h^0(A,L^{\otimes n})=n^g h^0(A,L)$.
A spanning family of that cardinality contributes
$g\log n+\log h^0(A,L)$ to $\log r$; the principal case has
$h^0(A,L)=1$.

\paragraph{An economical theta power and the arithmetic input.}
Suppose $L$ gives a principal polarization. Assume that for every
integer $m\ge1$, after a finite number field extension, the prescribed
construction gives $m^{2g}$ integral sections spanning $L^{\otimes m^2}$
with the same bounds $M_\sigma$, pulled back under the embeddings.
Assume the integral tangent lattice and its metric pull back
isometrically, and that for fixed $A_0,B_0\ge0$,
\[
 U\le h/2+A_0g+B_0g\log Y,\qquad Y=2+g+h,\qquad h\ge0.
\]
This includes the uniform finite-prime integrality of those sections;
it is not implied by their complex translation identities.

Put
\[
 Z=1+g+h+2A_0g+2B_0g\log Y,\qquad
 a=\lceil\sqrt Z\rceil,\quad m=3a,\quad k=3a^2,\quad n=9a^2 .
\]
Thus $n=m^2=3k$, $3Z\le k\le12Z$ and $n\le36Z$.
The tensor exponent is $n$; the torsion order is $m$.
With $r=n^g$, Proposition~\ref{geo-theta-jet-bound} yields
\[
 -\operatorname{av}_K\log\ell_\sigma
 \le\log k+\tfrac12\log(3e)+\frac{g\log(3k)+2U}{k}
 \le\frac43\log(12Z)+\frac56\log3+\frac12 .
\]
For the last inequality, use $\log(3k)>1$ and
$g+2U_+\le Z$. If $C_0=1+2A_0+2B_0$, then
$Z\le C_0Y\log Y$ and $\log Y\ge1$. Consequently
\begin{equation}\label{geo-logarithmic-tangent-line}
 -\operatorname{av}_K\log\|x\|
 \le C_{\rm line}\log(2+g+h),\qquad
 C_{\rm line}=\frac83+\frac43\log(12C_0)+\frac56\log3+\frac12 .
\end{equation}

The passage to degrees is supplied by the classical absolute Siegel
lemma. For a nonzero Hermitian sublattice $F$ of rank $r$, let $h_F$
be its logarithmic adelic height, using the given archimedean norms
and the finite-place unit balls $F_v$. Gaudron--R\'emond give
\[
 \inf_{0\ne y\in F_K\otimes_K\overline K}h_F(y)
       \le-\widehat\mu(F)+c_r,\qquad
 c_r=\frac12\sum_{j=2}^r\frac1j\le\frac12\log r .
\]
These are the same ordinary normalized degree and norm conventions
\cite[Sections~2.1--2.3 and Lemma~3.2]{gaudron-remond-minima}.

For completeness, the algebraic vector can be made integral without
changing its height. Choose $y$ within $\epsilon>0$ of the displayed
infimum, over a finite extension $K'$. Its saturated integral line
is $P=Iy$ for a fractional ideal $I$ of $\mathcal O_{K'}$.
Choose $a\ge1$ with $I^a=(b)$ and adjoin $\alpha$ with $\alpha^a=b$.
Fractional ideals have no torsion, so $I\mathcal O_L=(\alpha)$ in
$L=K'(\alpha)$. The vector $x=\alpha y$ is an integral generator
of $P\otimes\mathcal O_L$. The quotient by $P$ is projective,
so saturation survives base change and every finite gauge norm
of $x$ is one. The product formula now gives exactly
\[
 \operatorname{av}_{L}\log\|x\|=h_F(x)=h_F(y),\qquad
 \widehat\mu(F)-c_r-\epsilon
                  \le-\operatorname{av}_{L}\log\|x\| .
\]
No metric has changed.

Suppose the theta construction and \eqref{geo-logarithmic-tangent-line}
apply after every finite extension, with unchanged constants, rank,
and $h=(-d)_+$. Apply that line bound to $x$ and let
$\epsilon\downarrow0$. Since $r\le g$,
\[
 \widehat\mu(F)
 \le C_{\rm line}\log(2+g+(-d)_+)+c_r
 \le(C_{\rm line}+\tfrac12)\log(2+g+(-d)_+).
\]
The extension may depend on $\epsilon$; no bound on its degree is
needed. Multiplying by the actual selected ranks and summing supplies
Proposition~\ref{geo-selected-tangent-sum} with
$C_s=C_{\rm line}+\tfrac12$.
If the ambient abelian tangent rank is zero, the selected sum is zero
and none of these theta or Siegel inputs is needed.

The arithmetic inputs isolated by this argument are the primitive theta
section with its mean norm bound and integral normalized torsion
translates with the same supremum bounds at the chosen levels,
compatible with the prescribed field extensions and metric. The following
classical construction supplies them at positive even levels, which
suffice for the parameter choice.
At a degeneration the final valuation inequality is the elementary
identity
\[
 Q(\chi)+\frac{b_Q(y,\chi)}N+\frac{Q(y)}{N^2}
             =\frac{Q(N\chi+y)}{N^2}\ge0
\]
for a nonnegative quadratic form $Q$ and $N\ge1$. Applying it to
theta values still requires the coefficient valuation bound and the
identification of the prescribed rigidified torsion frame.
